\subsubsection{Rancangan Pelatihan Model}

Pelatihan dilakukan menggunakan YOLO11-seg dengan bobot pralatih sebagai titik awal. Bobot tersebut selanjutnya disesuaikan terhadap sembilan kelas makanan pada dataset penelitian sehingga proses pelatihan tidak dimulai dari bobot acak.

Data pelatihan diberikan kepada model dalam bentuk citra beserta \textit{ground truth} yang memuat label kelas dan poligon segmentasi setiap instans. Pada saat data dimuat, citra disesuaikan dengan ukuran masukan model dan dapat menerima augmentasi sesuai konfigurasi pelatihan.

Selama pelatihan, model menghasilkan prediksi yang dibandingkan dengan \textit{ground truth} untuk memperoleh nilai \textit{loss}. Nilai tersebut digunakan dalam proses \textit{backpropagation} dan pembaruan bobot. Proses ini dilakukan berulang selama sejumlah \textit{epoch} yang ditetapkan pada tahap implementasi.

Data validasi digunakan selama pelatihan untuk memantau perkembangan performa model tanpa digunakan untuk memperbarui bobot. Hasil validasi digunakan sebagai dasar pemilihan bobot model yang selanjutnya digunakan pada evaluasi akhir.

Model terpilih kemudian dievaluasi menggunakan data pengujian yang tidak digunakan untuk pembaruan maupun pemilihan bobot. Evaluasi dilakukan terhadap kemampuan model dalam mendeteksi, mengklasifikasikan, membentuk \textit{segmentation mask}, serta menghasilkan jumlah instans yang sesuai dengan \textit{ground truth}. Varian YOLO11-seg, ukuran masukan, jumlah \textit{epoch}, \textit{batch size}, dan konfigurasi pelatihan lainnya dicatat sesuai implementasi yang benar-benar digunakan.
